% =====================================================================
% Preámbulo — Informe Entrega 1, Grupo 4, Algoritmos Avanzados (UNSAAC)
% =====================================================================
\usepackage[utf8]{inputenc}
\usepackage[T1]{fontenc}
\usepackage[spanish,es-noquoting,es-tabla,es-nodecimaldot]{babel}
\usepackage{lmodern}
\usepackage{microtype}
\usepackage[a4paper,margin=2.5cm,headheight=22pt]{geometry}

\usepackage{amsmath,amssymb}
\usepackage{graphicx}
\usepackage[table]{xcolor}
\usepackage{booktabs,makecell,longtable,tabularx,multirow,array}
\usepackage{float}
\usepackage[font=small,labelfont={bf,color=guinda},format=hang]{caption}
\usepackage{enumitem}
\usepackage{fancyhdr}
\usepackage{titlesec}
\usepackage{listings}
\usepackage{algorithm}
\usepackage[noend]{algpseudocode}
\usepackage{tcolorbox}
\usepackage[hidelinks]{hyperref}

% ---------------------------------------------------------------- colores
\definecolor{guinda}{HTML}{6E0F1D}
\definecolor{dorado}{HTML}{B8860B}
\definecolor{doradoclaro}{HTML}{F6EFE3}
\definecolor{grisclaro}{HTML}{F3F3F3}

% ---------------------------------------------------------------- títulos
\titleformat{\section}{\Large\bfseries\color{guinda}}{\thesection.}{0.6em}{}[{\color{dorado}\titlerule[0.8pt]}]
\titleformat{\subsection}{\large\bfseries\color{guinda}}{\thesubsection}{0.6em}{}
\titleformat{\subsubsection}{\normalsize\bfseries\color{dorado!70!black}}{\thesubsubsection}{0.6em}{}
\setcounter{secnumdepth}{3}
\setcounter{tocdepth}{2}

% ---------------------------------------------------------------- encabezado
\pagestyle{fancy}
\fancyhf{}
\fancyhead[R]{\small\color{gray}UNSAAC · Algoritmos Avanzados · Grupo 4 · Entrega 1}
\fancyfoot[C]{\small\thepage}
\renewcommand{\headrulewidth}{0.4pt}
\renewcommand{\headrule}{\hbox to\headwidth{\color{dorado}\leaders\hrule height \headrulewidth\hfill}}

% ---------------------------------------------------------------- texto
\setlength{\parskip}{4pt}
\setlength{\parindent}{0pt}
\setlist{itemsep=2pt,topsep=3pt}
\renewcommand{\arraystretch}{1.15}
\newcolumntype{L}[1]{>{\raggedright\arraybackslash}p{#1}}
\newcolumntype{Y}{>{\raggedright\arraybackslash}X}

% ---------------------------------------------------------------- código
\lstset{
  basicstyle=\ttfamily\footnotesize,
  backgroundcolor=\color{grisclaro},
  frame=single, rulecolor=\color{grisclaro},
  breaklines=true, columns=fullflexible, keepspaces=true,
  literate={á}{{\'a}}1 {é}{{\'e}}1 {í}{{\'i}}1 {ó}{{\'o}}1 {ú}{{\'u}}1 {ñ}{{\~n}}1
           {Á}{{\'A}}1 {É}{{\'E}}1 {Í}{{\'I}}1 {Ó}{{\'O}}1 {Ú}{{\'U}}1
           {·}{{$\cdot$}}1 {→}{{$\to$}}1 {├}{{|}}1 {└}{{`}}1 {─}{{-}}1 {│}{{|}}1
}

% ---------------------------------------------------------------- pseudocódigo
\floatname{algorithm}{Algoritmo}
\renewcommand{\listalgorithmname}{Lista de algoritmos}
\algrenewcommand\algorithmicrequire{\textbf{Entrada:}}
\algrenewcommand\algorithmicensure{\textbf{Salida:}}
\algrenewcommand\algorithmicfunction{\textbf{función}}
\algrenewcommand\algorithmicif{\textbf{si}}
\algrenewcommand\algorithmicthen{\textbf{entonces}}
\algrenewcommand\algorithmicelse{\textbf{si no}}
\algrenewcommand\algorithmicfor{\textbf{para}}
\algrenewcommand\algorithmicforall{\textbf{para cada}}
\algrenewcommand\algorithmicdo{\textbf{hacer}}
\algrenewcommand\algorithmicreturn{\textbf{devolver}}
\algrenewcommand\algorithmicwhile{\textbf{mientras}}
\algnewcommand\Break{\textbf{romper}}

% ---------------------------------------------------------------- notas
\newtcolorbox{nota}[1][]{colback=doradoclaro,colframe=dorado,boxrule=0pt,leftrule=3pt,
  arc=0pt,left=6pt,right=6pt,top=4pt,bottom=4pt,fontupper=\small,#1}

% ---------------------------------------------------------------- capturas
% \captura[ancho]{archivo}{pie de figura}{etiqueta}
% Inserta una captura de pantalla de capturas/ con borde fino.
\newcommand{\captura}[4][0.92\textwidth]{%
  \begin{figure}[H]
  \centering
  {\setlength{\fboxsep}{0pt}\fcolorbox{gray!50}{white}{%
    \includegraphics[width=#1,height=9cm,keepaspectratio]{capturas/#2}}}
  \caption{#3}\label{#4}
  \end{figure}}

\newcommand{\Lmax}{L_{\max}}
\newcommand{\pbar}{\bar{p}}
